\chapter{Heat tensors and contact terms}
\label{mht:ch:heat-contacts}

Classical transfer applies a homotopy to smooth fields. Quantum
contraction instead evaluates two input covectors on a tensor of
fields, so it requires control of both factors at once. The paired
modes of Chapter~\ref{mht:ch:angular-cutoffs} provide this control at
positive heat time. They give a free BV complex and explicit
contractions of jets at the pole. Their generator also identifies
the limitation of this construction: it contains an angular
projection and is therefore nonlocal on the curve.

The tensor contracts all surface modes and evaluates jets at the
pole. Products of coincident kernels and collisions of poles
require further choices. A subtraction prescription determines
finite contacts, and convergence of a polar trace need not
preserve the boundary product.
The differential on supported normal jets determines which finite
contacts survive as classes. At a matrix background, a separate
Hilbert-scale estimate controls the inverse after the angular cutoff
is removed.

\section{The paired heat tensor}

Retain $E,d_E,\tau_E,h,V_\pm,\mathcal C,\mathcal R$ and the round
probability measure $d\mu$ from the preceding chapter. Choose an
orthonormal spherical basis $e_{\ell m}$ of $V_-$, where
$\ell\geq1$, $0\leq m\leq\ell$, and $e_{\ell m}$ has rotation
weight $-m$. Put
\[
 \lambda_\ell=\ell(\ell+1),\qquad
 f_{\ell m}=\mathcal Ce_{\ell m}.
\]
The polarization identity gives
\[
 \omega(e_{\ell m},f_{\ell'm'})
   =\kappa\delta_{\ell\ell'}\delta_{mm'},\qquad
 \omega(e,e')=\omega(f,f')=0.
\]
Cyclic self-adjointness therefore requires equal heat weights
on the two members of each pair.

Define $R_t$ to fix $1,\eta_0$ and to multiply
$e_{\ell m},f_{\ell m},d_Ee_{\ell m},d_Ef_{\ell m}$ by
$e^{-t\lambda_\ell}$. If $\Delta_-$ is the round scalar Laplacian
on $V_-$, its restriction to $V=\ker p_0$ is
\begin{equation}\label{mht:eq:paired-heat-operator}
 R_t|_V=e^{-t\Delta_-}\mathcal P_-
       +\mathcal Ce^{-t\Delta_-}\mathcal R\mathcal P_+.
\end{equation}
On functions include the residual term $p_0(f)1$, and on one-forms
set $R_t\eta=\eta_0p_1(\eta)+d_ER_th\eta$.

The signs of the tensor can be fixed on one pair. Write
$\langle x,y\rangle=\tau_E(xy)$. For even vectors $a,b$ and odd
vectors $A=d_Ea,B=d_Eb$, suppose
$\langle a,B\rangle=\kappa$ and
$\langle b,A\rangle=-\kappa$. Define
\begin{align}
 P^{\mathrm{pair}}&=\kappa^{-1}(a\otimes b-b\otimes a),\notag\\
 K^{\mathrm{pair}}=d_EP^{\mathrm{pair}}
   &=\kappa^{-1}(a\otimes B-B\otimes a-b\otimes A+A\otimes b).
 \label{mht:eq:heat-pair-signs}
\end{align}
The differential on a tensor has its Koszul sign.

In the convention
\[
 R_Kx=(-1)^{|x|+1}\langle x,K_{(1)}\rangle K_{(2)},
\]
$K^{\mathrm{pair}}$ contracts to the identity on these four vectors.
The residual tensor is
$K_{\mathrm{res}}=\kappa^{-1}(1\otimes\eta_0-\eta_0\otimes1)$.

For $t>0$ define
\begin{align}
 \mathcal P_t&=\frac1\kappa\sum_{\ell\geq1}\sum_{m=0}^{\ell}
   e^{-t\lambda_\ell}
   (e_{\ell m}\otimes f_{\ell m}-f_{\ell m}\otimes e_{\ell m}),
   \notag\\
 \mathcal K_t&=K_{\mathrm{res}}+d_E\mathcal P_t.
 \label{mht:eq:paired-heat-tensors}
\end{align}
For $0<\epsilon<T$, put
$P_{\epsilon,T}=\mathcal P_\epsilon-\mathcal P_T$.

\begin{theorem}\label{mht:thm:paired-heat}
The series in \eqref{mht:eq:paired-heat-tensors} converge with
all derivatives on the product sphere. Their limits lie in the
completed projective tensor product of smooth Neumann fields and
satisfy
\begin{equation}\label{mht:eq:coefficient-heat-ward}
 d_EP_{\epsilon,T}=\mathcal K_\epsilon-\mathcal K_T,
 \qquad P_{\epsilon,T}+P_{T,S}=P_{\epsilon,S}.
\end{equation}
Contraction with $\mathcal K_t$ is $R_t$. These operators form a
chain semigroup, preserve the unit, trace and Neumann condition,
and satisfy
\[
 \langle R_tx,y\rangle=\langle x,R_ty\rangle.
\]
They tend to the identity uniformly on bounded subsets of the
smooth field space.

For total derivative order $r\geq0$ and
$0<t\leq1$,
\begin{equation}\label{mht:eq:heat-tensor-bound}
 \|\mathcal K_t\|_{C^r}\leq C_rt^{-(r+7)/2}.
\end{equation}
The part with $\ell>L$ satisfies the same bound times
$e^{-t(L+1)^2/2}$, with a different constant.
\end{theorem}

\begin{proof}
An $L^2$-normalized spherical harmonic has $H^a$ norm bounded by
$C_a(1+\ell)^a$. Sobolev embedding therefore bounds its $C^r$
norm by $C_r(1+\ell)^{r+2}$. The estimate for $\mathcal C$ in
Proposition~\ref{mht:prop:angular-polarization} gives this same
safe bound for $f_{\ell m}$. One application of $d_E$ costs at
most one further power.

At total derivative order $r$, a summand of $\mathcal K_t$ is
bounded by $C_r(1+\ell)^{r+5}e^{-t\lambda_\ell}$. There are
$\ell+1$ summands at degree $\ell$. A Gaussian integral bounds
their sum by \eqref{mht:eq:heat-tensor-bound}. Splitting the
exponential into two equal factors gives the tail estimate.
The same argument applies to every derivative of $\mathcal P_t$.

Smooth functions and forms on a compact product identify with
the completed projective tensor products of their smooth spaces.
The Neumann jet constraints define closed complemented subspaces,
so they persist in these limits. The four-vector calculation
\eqref{mht:eq:heat-pair-signs} and the residual tensor now give
the asserted contraction operator.

All chain identities hold on each pair. Equal weights on paired
vectors give cyclic self-adjointness. Since $R_t1=1$, that identity
also gives $\tau_ER_t=\tau_E$. Finally, spectral convergence on
$V_-$, composed with the continuous maps $\mathcal C,\mathcal R$,
proves convergence in every smooth seminorm. The spectral tail
estimates are uniform when the input ranges over a bounded
smooth set.
\end{proof}

The tensor is smooth, but its contraction with an input still
uses $\tau_E$. Relative to an ordinary smooth density, the input
covectors retain the polar singularity. Already the residual
tensor distinguishes these two assertions about a kernel.

\section{The generator and curve locality}

Let $\Lambda_E=-\partial_tR_t|_{t=0}$. It acts by
$\lambda_\ell$ on each acyclic pair and by zero on $1,\eta_0$.
On scalars, write $r=|z|^2$ and
\[
 \Delta=-(1+r)^2\partial_z\partial_{\bar z}.
\]
This is the positive round Laplacian in our normalization.
Formula~\eqref{mht:eq:paired-heat-operator} determines
$\Lambda_E=\Delta\mathcal P_-+\mathcal C\Delta\mathcal R\mathcal P_+$
on $V$.

\begin{proposition}\label{mht:prop:heat-generator}
On the smooth scalar Neumann domain,
\begin{equation}\label{mht:eq:heat-generator}
 \Lambda_Ef=\Delta f+
   \frac{2(1+|z|^2)}{z}\partial_{\bar z}\mathcal P_+f.
\end{equation}
The second term extends smoothly across the marked point.
\end{proposition}

\begin{proof}
For a positive weight $m$, write $f=z^mF(r)$ and
$V(r)=-(1+r)^2F'(r)$. Then
$\mathcal Rf=\bar z^{m-1}\overline{V(r)}$. The equation
$d_E\mathcal Cv=\eta_0\overline v$ gives
$\mathcal C\Delta\mathcal Rf=z^mG(r)$, with
\[
 G'=rV''+mV',\qquad G=rV'+(m-1)V
\]
up to a holomorphic term. Substitution gives
$z^mG=\Delta(z^mF)+2(1+r)z^mF'$.

For $m=1$, the Neumann condition is $F(0)=0$, and the two
first-jet contributions in this expression cancel. For $m>1$
the first jet vanishes by weight. In the chart $w=1/z$, the
extra differential expression is
$-2(1+|w|^2)\bar w\partial_{\bar w}$, so it is smooth there.
The possible holomorphic difference is a global positive-weight
holomorphic function, and hence zero. Nonpositive weights give
$\Delta$, and constants give zero. Smooth Fourier convergence
and the estimates for $\mathcal C,\mathcal R$ extend the identity
to every smooth Neumann input.
\end{proof}

The projection $\mathcal P_+$ prevents locality. To see its
effect between disjoint supports, put $z=\rho e^{i\theta}$.
Choose an even nonnegative smooth function $b$, supported in a
small arc about zero, with mean
$\bar b=(2\pi)^{-1}\int b\,d\theta>0$. Choose an annular
function $a$ with $a(\rho_0)=0$ and $a'(\rho_0)=1$, and put
$f(\rho,\theta)=a(\rho)b(\theta)$. This function vanishes near
the opposite arc and near both poles.

\begin{proposition}\label{mht:prop:heat-nonlocal}
For this function,
\begin{equation}\label{mht:eq:heat-disjoint-response}
 (\Lambda_Ef)(\rho_0,\pi)
   =-\frac{1+\rho_0^2}{2\rho_0}\bar b,
 \qquad
 (R_tf)(\rho_0,\pi)
   =t\frac{1+\rho_0^2}{2\rho_0}\bar b+O(t^2).
\end{equation}
No finite-rank linear correction to $\Lambda_E$ makes it
support-preserving.
\end{proposition}

\begin{proof}
Off the angular diagonal, the kernel of the positive Fourier
projection is
\[
 k_+(\theta-\varphi)=\frac1{2\pi}
       \frac{e^{i(\theta-\varphi)}}{1-e^{i(\theta-\varphi)}}.
\]
Its real part is $-1/(4\pi)$. At $\theta=\pi$ its imaginary
part is odd in $\varphi$, so
$(\mathcal P_+b)(\pi)=-\bar b/2$. The local Laplacian vanishes
on the opposite arc.

The extra term in
\eqref{mht:eq:heat-generator} is
\[
 \frac{1+\rho^2}{\rho}
   \left(\partial_\rho+\frac{i}{\rho}\partial_\theta\right)
   \mathcal P_+f.
\]
At $\rho_0$, only the radial derivative remains. This proves
the first formula. On smooth inputs the spectral series is
twice differentiable at zero, giving the second.

Repeat the construction in arbitrarily many disjoint radial
annuli. On the opposite arc the resulting outputs are nonzero
and have disjoint radial supports, so they are linearly
independent. A correction restoring support preservation would
have to cancel all of them there. Its rank cannot be finite.
\end{proof}

A local cyclic operator exists if uniform ellipticity at the
pole is relinquished. Consider a first-order operator on
one-forms,
\[
 q(\eta\,d\bar z)=A\partial_z\eta+B\eta,
\]
with smooth coefficients, and set $q(E^0)=0$. Cyclicity means
$\tau_E((q\eta)\zeta)+\tau_E(\eta q\zeta)=0$.

\begin{proposition}\label{mht:prop:local-polar-gauge}
Cyclicity forces
\begin{equation}\label{mht:eq:local-cyclic-coefficients}
 2B=\partial_zA-\frac{2A}{z}.
\end{equation}
In particular $A(0)=0$. A real coefficient $A=-a$, with $a>0$
off the pole, therefore cannot give a uniformly elliptic
Dolbeault--Laplace principal part there.
The operator
\begin{equation}\label{mht:eq:local-polar-gauge}
 q_{\mathrm{loc}}(\eta\,d\bar z)
   =-r(1+r)\partial_z\eta+\tfrac12\bar z\eta
\end{equation}
is globally smooth, maps $E^1$ to $E^0$, is cyclic, and has
square zero. Its commutator
$\Lambda_{\mathrm{loc}}=[d_E,q_{\mathrm{loc}}]$ commutes with $d_E$.
\end{proposition}

\begin{proof}
Integration by parts on tests away from the pole gives
$2z^{-2}B-\partial_z(z^{-2}A)=0$. Thus $A/z$ extends smoothly,
proving \eqref{mht:eq:local-cyclic-coefficients} and $A(0)=0$.
The coefficients in \eqref{mht:eq:local-polar-gauge} satisfy this
identity. Its output and first holomorphic jet vanish at zero.

In the coordinate $w=1/z$, its action on $\eta_wd\bar w$ is
\[
 -(1+|w|^2)\partial_w\eta_w-\tfrac12\bar w\eta_w,
\]
which proves global smoothness.

Since $r(1+r)/z^2$ is bounded
near zero, the boundary term in polar integration by parts
tends to zero with the circle length. This proves cyclicity
on all smooth inputs. The square vanishes by degree, and the
commutation assertion follows from $d_E^2=0$.
\end{proof}

The degeneracy opens an infinite radial end for a natural
weighted realization. On radial scalars $f=F(r)$, put
$s=-\log r$ and $g(s)=F(e^{-s})$. Direct differentiation gives
\begin{equation}\label{mht:eq:local-radial-operator}
 \Lambda_{\mathrm{loc}}g=-(1+e^{-s})g''-\tfrac12g'.
\end{equation}

Fix $s_0\geq0$ and set
$a(s)=1+e^{-s}$, $w(s)=e^{s/2}/\sqrt{1+e^{-s}}$.
Since $(aw)'=w/2$, the expression is
$-w^{-1}\partial_s(aw\partial_s)$. Let $L$ be the Friedrichs
operator on $L^2([s_0,\infty),w\,ds)$ obtained by closing
$\int aw|g'|^2ds$ on $C_c^\infty((s_0,\infty))$.
This specifies a Dirichlet condition at $s_0$.

\begin{proposition}\label{mht:prop:local-radial-spectrum}
For this realization,
$\sigma_{\mathrm{ess}}(L)=[1/16,\infty)$. Its resolvent is
not compact, and $e^{-tL}$ is not trace class for any $t>0$.
\end{proposition}

\begin{proof}
The unitary substitution $u=w^{1/2}g$ gives
\[
 \widehat L=-\partial_s(a\partial_s)+V(s),\qquad
 V(s)=\frac{1-4e^{-s}-4e^{-2s}}{16(1+e^{-s})}.
\]
The coefficients approach those of $-\partial_s^2+1/16$
exponentially, with all derivatives. For $\xi\in\mathbb R$
and $\chi\in C_c^\infty(0,1)$ of norm one, the unit vectors
\[
 u_R(s)=R^{-1/2}\chi((s-R^2)/R)e^{i\xi s}
\]
escape every compact interval and satisfy
$\|(\widehat L-\xi^2-1/16)u_R\|=O(R^{-1})+O(e^{-R^2})$.
Weyl's criterion includes $[1/16,\infty)$ in the essential
spectrum.

For any $\varepsilon>0$, the quadratic form on a sufficiently
distant tail is bounded below by
$(1/16-\varepsilon)\|u\|^2$. Dirichlet--Neumann bracketing
separates that tail from a compact interval, whose resolvent
is compact. Hence no essential spectrum lies below
$1/16-\varepsilon$. Letting $\varepsilon$ decrease to zero
proves equality. This nonempty essential spectrum excludes
a compact resolvent and a compact, hence trace-class, heat
operator.
\end{proof}

After translating the endpoint to zero, the limiting operator
$-\partial_s^2+1/16$ has Dirichlet heat kernel
\[
 \frac{e^{-t/16}}{\sqrt{4\pi t}}
 \left(e^{-(s-s')^2/(4t)}-e^{-(s+s')^2/(4t)}\right).
\]
Odd reflection of the line Gaussian proves this formula.
It describes the limiting end, while
\eqref{mht:eq:local-radial-operator} retains variable coefficients.
The weighted Hilbert domain is additional analytic data. Its
spectrum makes no assertion about the cohomology of the smooth
Neumann complex.

\section{All surface modes and the free BV complex}

Return to the paired spectral regulator. Define
$Q_E^{\mathrm{GF}}=\Lambda_Eh$. The contraction identities give
\[
 [d_E,Q_E^{\mathrm{GF}}]=\Lambda_E,
 \qquad (Q_E^{\mathrm{GF}})^2=0.
\]
Both operators are continuous on smooth fields.

Let $U$ be a
closed compact oriented Riemannian surface and use the full space
\begin{equation}\label{mht:eq:heat-field-space}
 \mathcal F=
  \bigl(\Omega^\bullet(U)\widehat\otimes E
                   \otimes\mathfrak{gl}_N\bigr)[1].
\end{equation}
The tensor product is completed projectively. Surface forms
precede coefficient forms, and tensor operators use the Koszul
rule. Integration over $U$, the trace $\tau_E$, and
$\operatorname{tr}(XY)$ give a pairing of degree $-3$ before
the shift, and the odd BV pairing after it.

Set
\[
 Q=d_U+d_E,
 \qquad Q^{\mathrm{GF}}=d_U^*+Q_E^{\mathrm{GF}},
 \qquad [Q,Q^{\mathrm{GF}}]=\Delta_U+\Lambda_E.
\]
The cross commutators vanish by the tensor signs.

Let
$\mathbf K_t$ be the de Rham heat tensor on $U$, tensored
with $\mathcal K_t$ and the matrix coevaluation, with the
suspension signs induced by the pairing. Thus every smooth
surface mode is retained. Put
\begin{equation}\label{mht:eq:surface-heat-propagator}
 \mathbf P_{\epsilon,T}
   =\int_\epsilon^T(Q^{\mathrm{GF}}\otimes1)\mathbf K_t\,dt.
\end{equation}
For $0<\epsilon<T$ this is smooth, and
\begin{equation}\label{mht:eq:surface-heat-ward}
 Q\mathbf P_{\epsilon,T}=\mathbf K_\epsilon-\mathbf K_T.
\end{equation}
Indeed $Q\mathbf K_t=0$, and the commutator in the integral
acts as $-\partial_t$ on the heat semigroup.

Write $\mathcal O_n(\mathcal F)$ for continuous homogeneous
polynomials of degree $n$, defined by continuous graded symmetric
multilinear forms, and set
$\mathcal O_{\mathrm{poly}}=\bigoplus_{n\geq0}\mathcal O_n$.
Give $\mathcal O_n$ the strong dual topology of the completed
projective graded symmetric tensor power
$\widehat{\operatorname{Sym}}{}^n_\pi\mathcal F$:
its seminorms are uniform bounds on bounded sets of that space.
Use the product topology in field degree and formal completion
in $\hbar$; continuity in $\hbar$ is understood coefficientwise.
Contraction of two inputs with $\mathbf K_t$ defines an odd
constant second-order operator $\Delta_t$.
Contraction with $\mathbf P_{\epsilon,T}$ defines the even
operator $\partial_{\mathbf P_{\epsilon,T}}$, with convention
\[
 [Q,\partial_{\mathbf P_{\epsilon,T}}]
      =\Delta_\epsilon-\Delta_T.
\]
The derived bracket is
\[
 \{F,G\}_t=(-1)^{|F|}
  \bigl(\Delta_t(FG)-(\Delta_tF)G
                   -(-1)^{|F|}F\Delta_tG\bigr).
\]

\begin{proposition}\label{mht:prop:free-heat-bv}
On $\mathcal O_{\mathrm{poly}}(\mathcal F)[[\hbar]]$,
\begin{equation}\label{mht:eq:free-heat-flow}
 D_t=Q+\hbar\Delta_t,\qquad D_t^2=0,
 \qquad
 D_Te^{\hbar\partial_{\mathbf P_{\epsilon,T}}}
   =e^{\hbar\partial_{\mathbf P_{\epsilon,T}}}D_\epsilon.
\end{equation}
These formulas also define continuous coefficientwise operations
on $\bigl(\prod_{n\geq0}\mathcal O_n(\mathcal F)\bigr)[[\hbar]]$.
\end{proposition}

\begin{proof}
The smooth projective tensor can be evaluated on every
continuous multilinear coefficient, including marked jets.
The identity $Q\mathbf K_t=0$ gives $[Q,\Delta_t]=0$.
An odd constant contraction has square zero, and constant
contraction operators graded commute. Hence $D_t^2=0$.
Conjugating $D_T$ by the even exponential adds precisely
$\hbar(\Delta_\epsilon-\Delta_T)$, proving the intertwiner.

For output degree $n$ and coefficient $\hbar^r$, its
exponential uses only input degrees $n+2j$ with $0\leq j\leq r$.
Insertion of a fixed smooth tensor is a continuous linear map
between the completed symmetric tensor powers and preserves
bounded sets. Its transpose is therefore continuous for their
strong dual topologies. Each output coefficient is a finite sum
of these continuous contractions. This proves the assertion
about the product completion without summing infinitely many
contractions at one power of $\hbar$.
\end{proof}

At positive lower scale every fixed finite graph with continuous
polar interaction vertices is defined by these tensors. A local
counterterm theorem at zero scale requires a further argument:
Proposition~\ref{mht:prop:heat-nonlocal} excludes an appeal to
curve-local elliptic heat asymptotics for this generator.
If $U$ has boundary, the surface Stokes term also requires a
specified compatible BV--BFV domain.

\section{Pole jets and their subtractions}

The positive-scale tensor may be evaluated on the continuous
functionals
\[
 L_m(f)=\frac{\partial_{\bar z}^mf(0)}{m!},\qquad
 R_{m+1}(f)=\frac{\partial_z^{m+1}f(0)}{(m+1)!},\qquad m\geq1.
\]
They act on scalar acyclic fields and avoid the excluded first
holomorphic jet. Their oriented coefficient contraction is
$C_m(t)=(L_m\otimes R_{m+1})\mathcal P_t$.
After choosing a surface Gaussian polarization, the corresponding
component covariance is $\hbar C_m(t)$.

For an explicit calculation set
$x=(1-|z|^2)/(1+|z|^2)$ and let $P_\ell$ denote the Legendre
polynomial normalized by $P_\ell(1)=1$. Put
\[
 N_{\ell m}=\sqrt{\frac{(2\ell+1)(\ell-m)!}{(\ell+m)!}},
 \qquad
 e_{\ell m}=2^mN_{\ell m}
   \left(\frac{\bar z}{1+|z|^2}\right)^mP_\ell^{(m)}(x).
\]
Here $P_\ell^{(m)}$ means the ordinary $m$th derivative.
Legendre orthogonality gives norm one for $d\mu$.

The radial Dolbeault equation has solution
\begin{equation}\label{mht:eq:heat-radial-partner}
 \mathcal Ce_{\ell m}
   =\frac{N_{\ell m}z^{m+1}}2
       \int_{-1}^{x}(1+y)^mP_\ell^{(m)}(y)\,dy.
\end{equation}
Differentiating verifies $d_E\mathcal Ce=\eta_0\overline e$.
The integral vanishes to order at least $m+1$ at $x=-1$,
which gives smoothness at infinity. For $m=0$, its value at
$x=1$ is zero when $\ell\geq1$, giving the Neumann first jet.

\begin{lemma}\label{mht:lem:heat-legendre}
For $\ell,m\geq1$,
\begin{equation}\label{mht:eq:heat-legendre}
 \int_{-1}^{1}(1+x)^mP_\ell^{(m)}(x)\,dx
   =\frac2{(m-1)!}\prod_{j=1}^{m-1}
             \bigl(\ell(\ell+1)-j(j+1)\bigr).
\end{equation}
\end{lemma}

\begin{proof}
Integrate by parts $m$ times. All boundary terms at $-1$
vanish, and the remaining integral of $P_\ell$ is zero.
The derivatives at $1$ make the answer a polynomial in
$\lambda=\ell(\ell+1)$ of degree $m-1$, with leading
coefficient $2/(m-1)!$. For $\ell=1,\ldots,m-1$ the original
derivative vanishes. These roots determine the polynomial.
The empty product gives the case $m=1$.
\end{proof}

Let $\Theta(t)=\sum_{\ell\geq0}(2\ell+1)e^{-t\ell(\ell+1)}$.

\begin{proposition}\label{mht:prop:pole-heat-contacts}
For $t>0$ and $m\geq1$,
\begin{equation}\label{mht:eq:pole-heat-contacts}
 C_m(t)=\frac1{\kappa m!(m-1)!}
   \prod_{j=1}^{m-1}(-\partial_t-j(j+1))[\Theta(t)-1].
\end{equation}
In particular,
\begin{align}
 \kappa C_1(t)&=t^{-1}-\tfrac23+O(t),\notag\\
 \kappa C_2(t)&=\tfrac1{2t^2}-\tfrac1t+\tfrac{19}{30}+O(t),\notag\\
 \kappa C_3(t)&=\tfrac1{6t^3}-\tfrac2{3t^2}+\tfrac1t
                       -\tfrac{586}{945}+O(t).
 \label{mht:eq:first-pole-contacts}
\end{align}
For general $m$, the leading term is $1/(\kappa m!t^m)$.
\end{proposition}

\begin{proof}
The leading coefficients of $\bar z^m$ in $e_{\ell m}$ and
of $z^{m+1}$ in $\mathcal Ce_{\ell m}$ are
\[
 a_{\ell m}=N_{\ell m}\frac{(\ell+m)!}{m!(\ell-m)!},
 \qquad
 b_{\ell m}=\frac{N_{\ell m}}2
             \int_{-1}^{1}(1+x)^mP_\ell^{(m)}(x)\,dx.
\]
Their product is
$(2\ell+1)\prod_{j=1}^{m-1}(\lambda_\ell-j(j+1))/(m!(m-1)!)$.
The second orientation in $\mathcal P_t$ contributes zero by
weights. Sum the first against $\kappa^{-1}e^{-t\lambda_\ell}$.
Modes $1\leq\ell<m$ may be added because their coefficients
vanish, proving \eqref{mht:eq:pole-heat-contacts}.

Euler--Maclaurin summation gives
\[
 \Theta(t)\sim t^{-1}+\tfrac13+\tfrac{t}{15}
                   +\tfrac{4t^2}{315}+\tfrac{t^3}{315}
                   +\tfrac{4t^4}{3465}+\cdots.
\]
The integral of $(2x+1)e^{-tx(x+1)}$ on $[0,\infty)$ is
$1/t$. The endpoint half-value and the Bernoulli derivative
terms give the remaining coefficients. Its remainder estimate
also holds after any fixed number of $t$ derivatives.
Applying the differential polynomial gives
\eqref{mht:eq:first-pole-contacts}. Its highest derivative on
$t^{-1}$ gives the general leading term.
\end{proof}

For a finite set of pairs choose commuting formal sources
$J_m,K_m$ in the Gaussian polarization. Rotation weights kill
the pairings with distinct indices, so their generating series is
\[
 Z_t(J,K)=\exp\left(\hbar\sum_{m=1}^M C_m(t)J_mK_m\right).
\]
Let $\operatorname{PP}$ and $\operatorname{FP}$ denote the
negative-power part and constant part of the small-$t$ expansion.
Then, coefficientwise in sources and $\hbar$,
\[
 e^{-\hbar\sum_m\operatorname{PP}C_m(t)J_mK_m}Z_t
   \longrightarrow
 e^{\hbar\sum_m\operatorname{FP}C_m(t)J_mK_m}.
\]
This defines the Gaussian multipoint coefficients in the selected
component sector. Gauge invariance and comparison with sourced
closed observables require their own maps.

\begin{proposition}\label{mht:prop:heat-jet-subtraction}
On a finite jet system closed under the dual differential,
the divergent parts of the paired tensors obey
$\mathrm{CT}_K=d_E\mathrm{CT}_P$. Their subtraction preserves
the free Ward identity. That identity does not fix the finite
parts.
\end{proposition}

\begin{proof}
Choose $N\geq M+1$, take scalar jets through order $N$ and
one-form jets through order $N-1$, and impose the Neumann
relation on scalar jets. The dual of $d_E$ sends a one-form jet to a scalar jet of at
most order $N$. These tests therefore form a finite differential
system containing the selected pairs.
Evaluate $d_E\mathcal P_t=\mathcal K_t-K_{\mathrm{res}}$ on
it. Each asymptotic coefficient commutes with the finite
differential, proving the subtraction identity.

Replacing the heat parameter by $t+ct^2$ retains the paired
identity but changes the first covariance to
$\kappa^{-1}(t^{-1}-2/3-c+O(t))$. Its finite part varies
by $-c/\kappa$. This is a reparametrized prescription. It need
not be a semigroup in the new parameter.
\end{proof}

\section{Coincident kernels}

Linear approximation on smooth inputs does not by itself
control products of kernels. Consider the ordinary scalar
sharp projector on the sphere, independently of the paired
projection of Chapter~\ref{mht:ch:angular-cutoffs}. At the
north pole its kernel relative to $d\mu$ is
\[
 D_L(x)=\sum_{\ell=0}^L(2\ell+1)P_\ell(x),\qquad x=\cos\theta.
\]

\begin{proposition}\label{mht:prop:sharp-contact-obstruction}
No sequence of counterterms supported at the north pole makes
$D_L^2$ converge as a distribution. If $a$ is a nonzero
nonnegative smooth radial test supported in a compact annulus
in $0<\theta<\pi$, then
\begin{equation}\label{mht:eq:sharp-annular-contact}
 \lim_{L\to\infty}\frac1L\int a(\theta)D_L(\cos\theta)^2d\mu
   =\frac1{2\pi}\int\frac{a(\theta)}{\sin^2(\theta/2)}d\theta>0.
\end{equation}
\end{proposition}

\begin{proof}
Christoffel--Darboux gives
$(1-x)D_L(x)=(L+1)(P_L(x)-P_{L+1}(x))$.
Uniformly on a compact interior interval, Darboux's formula is
\[
 P_L(\cos\theta)=\sqrt{\frac2{\pi L\sin\theta}}
   \cos((L+\tfrac12)\theta-\tfrac\pi4)+O(L^{-3/2}).
\]
This is the zero-parameter Jacobi case of
\href{https://dlmf.nist.gov/18.15.E4_5}{DLMF 18.15.4\_5}.
Consequently the leading term of $D_L$ is
\[
 \sqrt{\frac{2L}{\pi\sin\theta}}
       \frac{\sin((L+1)\theta-\pi/4)}{\sin(\theta/2)}.
\]

Squaring and integrating against the radial measure
$\frac12\sin\theta\,d\theta$ gives
\eqref{mht:eq:sharp-annular-contact}. The oscillatory cosine
has vanishing integral in the limit. Every pole-supported
counterterm vanishes on $a$, while this positive integral
grows linearly.
\end{proof}

The same kernel has exact moments
\begin{align*}
 \tfrac12\int_{-1}^{1}D_L^2dx&=(L+1)^2,\\
 \tfrac12\int_{-1}^{1}(x-1)D_L^2dx&=-(L+1),\\
 \tfrac12\int_{-1}^{1}(1-x)^2D_L^2dx
   &=(L+1)^2\left(\frac1{2L+1}+\frac1{2L+3}\right).
\end{align*}
Orthogonality proves the first. The three-term recurrence for
$xP_\ell$ proves the second, and the squared
Christoffel--Darboux identity proves the third. These moments belong to the ordinary sharp kernel. A comparison
with the paired regulator would require its individual components.

There is also a finite comparison between different heat
prescriptions. In a separate complex normal plane, with real
coordinates $(x,y)$, put
\[
 \Delta_\perp=\partial_x^2+\partial_y^2,
 \qquad \delta_t(x,y)=\frac1{4\pi t}e^{-(x^2+y^2)/(4t)}.
\]
These are scalar coefficient densities after contraction.

\begin{proposition}\label{mht:prop:normal-heat-contact}
For $n\geq2$, the distributional finite part is
\begin{equation}\label{mht:eq:normal-heat-contact}
 \operatorname{FP}_{t=0}\delta_t^n
   =\frac{\Delta_\perp^{n-1}\delta_0}
          {(4\pi)^{n-1}n^n(n-1)!}.
\end{equation}
The joint cubic prescription minus the prescription which
first takes the binary finite part and then heat-regularizes
both remaining factors is
\begin{equation}\label{mht:eq:normal-triple-difference}
 \frac{37}{55296\pi^2}\Delta_\perp^2\delta_0.
\end{equation}
\end{proposition}

\begin{proof}
The identities
\[
 \delta_t^n=\frac{\delta_{t/n}}{n(4\pi t)^{n-1}},
 \qquad \delta_{t/n}=e^{(t/n)\Delta_\perp}\delta_0
\]
give \eqref{mht:eq:normal-heat-contact} by taking the constant
coefficient.
The binary and joint cubic finite parts are respectively
$\Delta_\perp\delta_0/(16\pi)$ and
$\Delta_\perp^2\delta_0/(864\pi^2)$.

For the successive prescription,
$\delta_t\Delta_\perp\delta_t=\frac12\partial_t(\delta_t^2)$,
whose finite part is $\Delta_\perp^2\delta_0/(128\pi)$.
Multiplication by the preceding binary coefficient $1/(16\pi)$
gives $\Delta_\perp^2\delta_0/(2048\pi^2)$. Subtracting yields
\eqref{mht:eq:normal-triple-difference}.
\end{proof}

The comparison is between two declared prescriptions on scalar
densities. It supplies a local contact between them. A product
of arbitrary singular Dolbeault currents, or coherent collision
operations in all arities, requires additional extension data.

\section{Two poles approaching one}

The spectral coordinate has its own contact, separate from
the normal plane. Work in a fixed chart and use compactly
supported smooth coefficients $\eta$ with measure $d^2z=dx\,dy$.
For $a\ne0$ define the locally integrable distribution
$K_a(z)=(z^2-a^2)^{-1}$ and functionals
\[
 \tau_a(\eta)=-\frac\kappa\pi\langle K_a,\eta\rangle,
 \qquad
 \tau_0(\eta)=-\frac\kappa\pi
       \langle\operatorname{PV}(z^{-2}),\eta\rangle.
\]
Here $\operatorname{PV}(z^{-2})=-\partial_z(1/z)$, and the
orientation fixes $\partial_{\bar z}(1/z)=\pi\delta_0$.
Thus $\tau_0(\eta)$ is the circular double-pole trace of
$\eta\,d\bar z$.

\begin{theorem}\label{mht:thm:two-pole-contact}
There is an exact distributional identity
\begin{equation}\label{mht:eq:two-pole-contact}
 K_a=\frac12\int_{-1}^{1}
       \operatorname{PV}\frac1{(z-ta)^2}\,dt
       -\frac{\pi\bar a}{2a}\int_{-1}^{1}\delta_{ta}\,dt.
\end{equation}
Along $a=\rho e^{i\theta}$ with $\rho\downarrow0$,
\[
 \tau_a=\tau_0+\kappa e^{-2i\theta}\operatorname{ev}_0
         +O(\rho^2).
\]
For tests in a fixed compact chart the error is bounded by
$C|a|^2\|\eta\|_{C^3}$, uniformly in $\theta$.
The corrected functional
\begin{equation}\label{mht:eq:segment-corrected-trace}
 \tau_a^{\mathrm{seg}}(\eta)
   =\tau_a(\eta)-\frac{\kappa\bar a}{2a}
                            \int_{-1}^{1}\eta(ta)\,dt
\end{equation}
converges to $\tau_0$ with the same bound and satisfies
\begin{equation}\label{mht:eq:segment-ward}
 \tau_a^{\mathrm{seg}}(\partial_{\bar z}f)
   =\frac\kappa2\int_{-1}^{1}\partial_zf(ta)\,dt
\end{equation}
for compactly supported smooth $f$.
\end{theorem}

\begin{proof}
Differentiate the translated locally integrable distribution:
\[
 \frac{d}{dt}\frac1{z-ta}
   =a\operatorname{PV}\frac1{(z-ta)^2}-\pi\bar a\delta_{ta}.
\]
Integration from $-1$ to $1$, followed by division by $2a$,
gives \eqref{mht:eq:two-pole-contact}. Translation is
continuous on distributions, so the first average tends to
$\operatorname{PV}(z^{-2})$ and the second gives the stated
directional contact.

On a fixed compact neighborhood, $1/z$ is integrable and
its derivative has order at most one. Symmetric Taylor
expansion of the translated test cancels the linear term.
The remainder in $C^1$ is bounded by
$C|a|^2\|\eta\|_{C^3}$. The symmetric delta average differs
from evaluation at zero by $C|a|^2\|\eta\|_{C^2}$.
These two estimates prove both convergence bounds.

The simple-pole distributional identity gives
$\tau_a(\partial_{\bar z}f)=\kappa(f(a)-f(-a))/(2a)$.
The chain rule along $ta$ separates this difference into
holomorphic and antiholomorphic derivatives. Subtracting
the segment term leaves \eqref{mht:eq:segment-ward}.
\end{proof}

For $f=\bar z$ near the segment, cut off outside a larger
disc, the raw expression is $\kappa\bar a/a$, while the
double-pole expression is zero. The contact is therefore
detected on smooth tests even though holomorphic residue
formulas alone do not see it.

A correction to the trace does not automatically specify an
algebra of boundary functions. Put
$\ell_a(f)=\frac12\int_{-1}^{1}\partial_zf(ta)dt$.
With the same cutoff convention,
\begin{equation}\label{mht:eq:segment-kernel-defect}
 \ell_a(\bar z)=\ell_a(z^2)=0,
 \qquad \ell_a(\bar z z^2)=\tfrac23|a|^2.
\end{equation}
Hence $\ker\ell_a$ is not a subalgebra when $a\ne0$.

In contrast, the raw condition $f(a)=f(-a)$ is multiplicative.
Its normalized difference tends, along the ray, to
$\partial_zf(0)+e^{-2i\theta}\partial_{\bar z}f(0)$.
This is a point derivation, so its kernel is an algebra, but
it imposes a directional boundary condition. Comparing it
with the holomorphic Neumann condition requires a map
preserving the boundary product as well as the corrected trace.

\section{Relative normal insertions and supported contacts}

The normal contact in \eqref{mht:eq:normal-triple-difference} is a
nonzero scalar distribution. To decide whether it represents an
obstruction, one must place it in a complex with its allowed
primitives. We first normalize its insertion on a disk, retaining the
outer boundary, and then compute the supported Dolbeault complex.

Write $p=x+iy$ for this normal coordinate, independently of the
spectral coordinate $z$. Orient $D_R=\{|p|\leq R\}$ by
$dx\wedge dy$ and its boundary $S_R$ by the induced orientation.
Normalize the area delta by $\langle\delta_0,\phi\rangle=\phi(0)$.
Then
\[
 \bar\partial(1/p)=\pi\delta_0\,d\bar p,\qquad
 d(dp/p)=2\pi i\delta_0\,dx\wedge dy.
\]

Let $\mathcal K^\bullet(D_R)$ be the de Rham currents that are
smooth on a fixed outer collar. Give them the strong distribution
topology in the interior and the smooth topology on that collar.
Restriction $i^*$ to $S_R$ is therefore continuous. The relative
complex and its charge are
\begin{align}
 \mathcal N_R^n
   &=\mathcal K^n(D_R)\oplus\Omega^{n-1}(S_R),\notag\\
 d_{\mathrm{rel}}(\omega,\gamma)
   &=(d\omega,i^*\omega-d\gamma),\notag\\
 \operatorname{Ch}(\omega,\gamma)
   &=\frac1{2\pi i}
       \left(\int_{D_R}\omega-\int_{S_R}\gamma\right)
       \quad(n=2).
       \label{mht:eq:relative-normal-charge}
\end{align}
The charge is zero in other degrees. Stokes' theorem makes it a
cochain map to $\mathbb C[-2]$, with its generator in degree two.

For normal heat time $t>0$, independently of the spectral heat time,
put
\begin{align*}
 Q_t&=1-e^{-R^2/(4t)},&
 \beta_t^H&=(1-e^{-|p|^2/(4t)})\frac{dp}{p},&
 \widetilde\beta_t^H&=Q_t^{-1}\beta_t^H,\\
 q_\epsilon&=\frac{R^2}{R^2+\epsilon^2},&
 \beta_\epsilon^R&=\frac{\bar p\,dp}{|p|^2+\epsilon^2},&
 \widetilde\beta_\epsilon^R&=q_\epsilon^{-1}\beta_\epsilon^R
       \quad(\epsilon>0).
\end{align*}
The superscripts distinguish heat and rational representatives.
Define $U_t^H=(d\widetilde\beta_t^H,0)$ and
$U_\epsilon^R=(d\widetilde\beta_\epsilon^R,0)$ in $\mathcal N_R^2$.

\begin{proposition}\label{mht:prop:relative-normal-insertion}
Both insertions have charge one and boundary primitive $i\,d\theta$.
Their difference is the relative boundary
\begin{equation}\label{mht:eq:normal-regulator-homotopy}
 U_t^H-U_\epsilon^R
   =d_{\mathrm{rel}}
        (\widetilde\beta_t^H-\widetilde\beta_\epsilon^R,0).
\end{equation}
Differences of normalized primitives give additive homotopies between
any two positive scales. Moreover,
$\widetilde\beta_t^H\to dp/p$ in the declared current and collar
topologies as $t\downarrow0$.
\end{proposition}

\begin{proof}
The function $(1-e^{-r^2/(4t)})/r^2$ is smooth at zero, so
$\beta_t^H$ is smooth there. Polar differentiation gives
\[
 d\beta_t^H=\frac{ir}{2t}e^{-r^2/(4t)}dr\wedge d\theta
          =2\pi i\delta_t\,dx\wedge dy.
\]
Both normalized primitives restrict to $i\,d\theta$ on $S_R$.
Stokes gives charge one, and equality of the boundary restrictions
proves \eqref{mht:eq:normal-regulator-homotopy}. Differences telescope,
which proves additivity.

Away from zero the heat primitive converges in every smooth seminorm.
Its local $L^1$ difference from $dp/p$ is bounded by a constant times
$\int_0^R e^{-r^2/(4t)}dr=O(\sqrt t)$.
This gives convergence uniformly on bounded test sets.
Since $Q_t\to1$, normalization preserves that convergence and the
smooth collar limit.
\end{proof}

This is a comparison of independently specified insertion maps.
It preserves relative charge and the stated support and boundary
conditions. It does not define a coincident product of singular currents.

\subsection{The differential on supported jets}

Set $\rho_t=Q_t^{-1}\delta_t$ on $D_R$ and
\[
 c_n=\frac1{(4\pi)^{n-1}n^n(n-1)!},\qquad n\geq2.
\]
The factor $Q_t^{-n}-1$ is exponentially small with all fixed
$t$ derivatives. The finite-part calculation of
Proposition~\ref{mht:prop:normal-heat-contact} therefore gives
$\operatorname{FP}_{t=0}\rho_t^n=c_n\Delta_\perp^{n-1}\delta_0$.
Its relative de Rham primitive is supported at the origin:
\begin{align}
 V_n&=c_n\bigl(
     \partial_x\Delta_\perp^{n-2}\delta_0\,dy
    -\partial_y\Delta_\perp^{n-2}\delta_0\,dx\bigr),\notag\\
 d_{\mathrm{rel}}(V_n,0)
   &=(c_n\Delta_\perp^{n-1}\delta_0\,dx\wedge dy,0).
       \label{mht:eq:real-contact-primitive}
\end{align}
Direct differentiation proves the identity, and support makes the
boundary restriction zero.

The Dolbeault quotient retains more information than this real charge.
Define the finite supported normal-jet complex by
\begin{align}
 \mathcal J_p^0&=\mathbb C[\xi,\eta]\delta_0,&
 \mathcal J_p^1&=\mathbb C[\xi,\eta]\delta_0\,d\bar p,\notag\\
 \xi&=\partial_p,&\eta&=\partial_{\bar p},\notag\\
 \bar\partial(\xi^a\eta^b\delta_0)
   &=\xi^a\eta^{b+1}\delta_0\,d\bar p.
       \label{mht:eq:supported-normal-jets}
\end{align}
The symbols denote distribution derivatives without factorial
normalization. Degree is Dolbeault degree. Use the inductive topology
of finite normal-order spaces, with smooth tangential coefficients at
fixed compact supports when these are present. Formal completions are
coefficientwise, with finite normal order at each coefficient.

\begin{theorem}\label{mht:thm:supported-normal-contraction}
The complex $\mathcal J_p$ contracts onto
$\mathbb C[\xi]\delta_0\,d\bar p$ in degree one. Let $\pi$ retain
only degree-one terms with antiholomorphic derivative order zero,
and let $\iota$ be their inclusion. The homotopy is zero in degree
zero and is
\[
 h_p(\xi^a\eta^b\delta_0\,d\bar p)=
 \begin{cases}
  \xi^a\eta^{b-1}\delta_0,&b\geq1,\\
  0,&b=0.
 \end{cases}
\]
It satisfies
\[
 \bar\partial h_p+h_p\bar\partial=1-\iota\pi,
 \qquad h_p^2=\pi h_p=h_p\iota=0.
\]
Thus $H^0(\mathcal J_p)=0$ and
$H^1(\mathcal J_p)=\mathbb C[\partial_p]\delta_0\,d\bar p$.
The maps are continuous, commute with holomorphic normal
differentiation, and extend after tensoring with an independent
cochain coefficient space in the stated finite-order topology.
\end{theorem}

\begin{proof}
On a degree-zero monomial $h_p\bar\partial$ is the identity.
On a degree-one monomial with $b\geq1$, $\bar\partial h_p$ is the
identity. On the remaining monomials both sides of $1-\iota\pi$
vanish. These cases also prove the side conditions.

The homotopy lowers normal order, so it is continuous at every
finite stage. The graded tensor rule cancels the two terms
involving an independent coefficient differential. This proves
the tensor assertion without passing to an unrestricted completion.
\end{proof}

In particular, $\delta_0\,d\bar p$ has no supported primitive.
Such a primitive would require negative antiholomorphic derivative
order. The higher heat contacts behave differently, since
$\Delta_\perp=4\partial_p\partial_{\bar p}$ gives
\begin{equation}\label{mht:eq:dolbeault-contact-primitive}
 \Delta_\perp^{n-1}\delta_0\,d\bar p
 =\bar\partial\bigl(
   4^{n-1}\partial_p^{n-1}\partial_{\bar p}^{n-2}\delta_0\bigr).
\end{equation}
The cubic difference in \eqref{mht:eq:normal-triple-difference}
therefore satisfies
\begin{equation}\label{mht:eq:cubic-contact-primitive}
 \frac{37}{55296\pi^2}\Delta_\perp^2\delta_0\,d\bar p
 =\bar\partial\left(
    \frac{37}{3456\pi^2}\partial_p^2\partial_{\bar p}\delta_0\right).
\end{equation}

The primitive requires a mixed normal jet in degree zero.
A pure holomorphic residue module does not contain the heat contact.
Adjoining the contact while excluding its displayed primitive
changes the complex and hence the obstruction question.

\subsection{Two normal directions and a surface boundary}

Tensor products below use independent normal variables, at finite
normal order and fixed compact test support. They do not multiply
arbitrary distributions in a single variable. Suppose $U_i,U_i'$
are closed cochains of degree $a_i$, and
$dH_i=U_i-U_i'$ with $|H_i|=a_i-1$. The two ordered homotopies are
\begin{align*}
 H_{12}&=H_1\otimes U_2+(-1)^{a_1}U_1'\otimes H_2,\\
 H_{21}&=(-1)^{a_1}U_1\otimes H_2+H_1\otimes U_2'.
\end{align*}

Both differentiate to $U_1\otimes U_2-U_1'\otimes U_2'$.
Their difference is
\begin{equation}\label{mht:eq:normal-corner-homotopy}
 H_{12}-H_{21}=(-1)^{a_1+1}d(H_1\otimes H_2).
\end{equation}
Indeed, the first difference expands as
$H_1\otimes dH_2-(-1)^{a_1}dH_1\otimes H_2$,
which is the displayed differential. The normal charges have
$a_i=2$. Adjoining an independent closed odd label to each factor,
placed last, changes $a_i$ to three and reverses the corner sign.

The tensor contraction for $\mathcal J_u\otimes\mathcal J_v$ is
\[
 h_u\otimes1+(\iota_u\pi_u)\otimes h_v.
\]
It is interpreted as a graded tensor of maps, so the second term
has a minus sign after a degree-one residue in the first factor.
The retained cohomology is
\[
 \mathbb C[\partial_u,\partial_v]
     \delta_u\delta_v\,d\bar u\,d\bar v
\]
in degree two. Expanding the tensor differential and using the two
one-variable contraction identities proves this assertion.

To retain contacts along the union of the two divisors, one must also
retain its faces. In one normal variable let $\mathcal L$ be the
Dolbeault coefficient complex localized at the divisor, with finite
pole order. Its principal-value section uses
\[
 E_m(p)=\frac{(-1)^{m-1}}{(m-1)!}
                \partial_p^{m-1}(1/p).
\]
Let $\mathcal E=\sigma(\mathcal L)+\mathcal J$ inside currents,
where $\sigma$ is this graded section and $\mathcal J$ is the full
supported jet complex. In the graded splitting,
\[
 \mathcal E=\mathcal L\oplus\mathcal J,\qquad
 d(a,s)=(d_{\mathcal L}a,d_{\mathcal J}s+\mathsf Ra),\qquad
 \mathsf R=d\sigma-\sigma d_{\mathcal L}.
\]

The identities $pE_1=1$ and $pE_m=E_{m-1}$ make the section
well-defined on localized coefficients. Restriction off the divisor
identifies the quotient with $\mathcal L$.
The equation $d_{\mathcal J}\mathsf R+\mathsf R d_{\mathcal L}=0$
follows from $d^2=0$.

For two variables the supported kernel of
$\mathcal E_u\boxtimes\mathcal E_v\to
 \mathcal L_u\boxtimes\mathcal L_v$ has graded space
\begin{equation}\label{mht:eq:normal-face-kernel}
 \mathcal J_{uv}=
 (\mathcal J_u\boxtimes\mathcal L_v)
 \oplus(\mathcal L_u\boxtimes\mathcal J_v)
 \oplus(\mathcal J_u\boxtimes\mathcal J_v).
\end{equation}
Its differential retains both face-to-corner maps. With
$r_1=\bar\partial(1/p)$, the signs are
\begin{align*}
 \bar\partial(E_1(u)E_1(v))
   &=r_1(u)E_1(v)+E_1(u)r_1(v),\\
 \bar\partial(r_1(u)E_1(v))&=-r_1(u)r_1(v),\\
 \bar\partial(E_1(u)r_1(v))&=+r_1(u)r_1(v).
\end{align*}
Deleting either face destroys the cancellation in the second
differential.

Now let $\mathcal B$ be an independently given cochain complex
with cochain map $\mathsf W:\mathcal B\to\mathcal L$.
The pullback of the extension is
\[
 \mathcal E_{\mathsf W}=\mathcal B\oplus\mathcal J,
 \qquad d(b,s)=(d_{\mathcal B}b,d_{\mathcal J}s+\mathsf R\mathsf Wb).
\]
It maps to $\mathcal E$ by $(b,s)\mapsto(\mathsf Wb,s)$.
Both the square-zero and chain-map identities follow from
$d_{\mathcal L}\mathsf W=\mathsf Wd_{\mathcal B}$ and the identity
for $\mathsf R$. This construction supplies a coefficient extension,
without a product of meromorphic functions with arbitrary supported
currents.

\begin{proposition}\label{mht:prop:relative-contact-totalization}
Suppose these complexes and maps are given on a surface $U$ and its
boundary $\Gamma$. Let the degree-zero cochain restrictions $\rho$
commute with $\mathsf W$ and satisfy
$\rho_{\mathcal J}\mathsf R_U=\mathsf R_\Gamma\rho_{\mathcal L}$.
The relative complex has differential
\begin{equation}\label{mht:eq:relative-contact-differential}
 D(b,s;c,v)=
 \left(
 \begin{array}{c}
 d_{\mathcal B}b,\quad d_{\mathcal J}s+\mathsf R_U\mathsf W_Ub;\\
 \rho_{\mathcal B}b-d_{\mathcal B}c,\quad
 \rho_{\mathcal J}s-d_{\mathcal J}v
                         -\mathsf R_\Gamma\mathsf W_\Gamma c
 \end{array}\right).
\end{equation}
Here $(b,s)$ has degree $n$ over $U$, and $(c,v)$ has degree
$n-1$ over $\Gamma$. One has $D^2=0$.
For two normal directions, $\mathcal J$ includes the full kernel
\eqref{mht:eq:normal-face-kernel}.
\end{proposition}

\begin{proof}
The restriction hypotheses make
$\mathcal E_{\mathsf W,U}\to\mathcal E_{\mathsf W,\Gamma}$ a
cochain map. For any cochain map $f:A\to B$, applying
$(a,b)\mapsto(d_Aa,fa-d_Bb)$ twice gives zero, since
$fd_A=d_Bf$. Substitution of the two extension differentials
gives \eqref{mht:eq:relative-contact-differential}.
\end{proof}

Thus a closed degree-one element satisfies all four equations
\[
 \begin{aligned}
 d_{\mathcal B}b&=0,&
 d_{\mathcal J}s+\mathsf R_U\mathsf W_Ub&=0,\\
 \rho_{\mathcal B}b-d_{\mathcal B}c&=0,&
 \rho_{\mathcal J}s-d_{\mathcal J}v
                        -\mathsf R_\Gamma\mathsf W_\Gamma c&=0.
 \end{aligned}
\]
If a sourced BV--BFV action realizes this diagram
in its admissible local cyclic counterterm complex, a degree-zero
counterterm changes its first quantum residual by $DX$.
For $X=(x,u;y,w)$ the boundary contact component is
\[
 \rho_{\mathcal J}u-d_{\mathcal J}w
                       -\mathsf R_\Gamma\mathsf W_\Gamma y.
\]
An admissible regulator transport giving such a boundary preserves
the residual class.

The primitive \eqref{mht:eq:cubic-contact-primitive}, tensored with
a closed independent coefficient, kills that supported coefficient
class. It does not kill additional field or boundary components.
Constructing the native sourced action, its local cyclic domain,
compatible restrictions and quantum measure is still required.
A compact-group disk preparation requires its own real integration
cycle, measure and sewing comparison. Its state is not determined
by the charge-one normal insertion.

\section{Uniform inverses at a matrix background}

The finite inverse in \eqref{mht:eq:angular-background-inverse} is
polynomial in $J$, but its degree grows with the cutoff. To take a
limit one needs a common bound. The adapted norm
\eqref{mht:eq:angular-adapted-norm} reveals an infinite square-zero
block whose removal would destroy the operator-norm approximation.

For $s\geq3$, let $H_s^-$ be the closure of $V_-$ in ordinary
round $H^s$, and complete $V_-\oplus V_+$ in the adapted norm
to obtain $\mathfrak V_s$. Its complex linear Hilbert coordinates are
\[
 \mathfrak V_s\longrightarrow H_s^-\oplus\overline{H_s^-},
 \qquad f\longmapsto(f_-,\overline{\mathcal Rf_+}).
\]
The second bar denotes the conjugate vector space, compensating for
the antilinearity of $\mathcal R$. The one-derivative gain of
$\mathcal C$ embeds $\mathfrak V_s$ continuously in ordinary $H^s$.
Its vectors satisfy the mean and Neumann constraints, and the
intersection of the scale is the original smooth space $V$.

Let $\mathcal P_0^{\mathrm{rot}}$ project to rotation weight zero,
distinct from the scalar residual mean $p_0$. Put
\[
 X=h\eta_0,\qquad
 \mathsf Nf=h(\eta_0\mathcal P_0^{\mathrm{rot}}f_-),\qquad
 K=X-\mathsf N.
\]

\begin{proposition}\label{mht:prop:staggered-nonsmoothing-block}
The operator $\mathsf N$ is bounded on $\mathfrak V_s$ and satisfies
$\mathsf N^2=0$. The remainder $K$ maps $\mathfrak V_s$ boundedly
to $\mathfrak V_{s+1}$. The actual staggered projections
$\widehat\Pi_L$ are orthogonal at every index, converge strongly
to the identity and commute with $\mathsf N$.
\end{proposition}

\begin{proof}
Write $f_-=r+w$, with $r$ of weight zero and $w$ of strictly
negative weights. Then
\[
 \mathsf Nr=\mathcal C\overline r,\qquad
 \mathcal R\mathsf Nr=\overline r.
\]
These bars are pointwise conjugation. Thus $\mathsf N$ sends the
radial summand isometrically to its positive partner in the adapted
norm. It kills positive weights, proving boundedness and square zero.

The vector $Xw$ has nonpositive weights and gains one ordinary
Sobolev derivative. The finite-rank corrections in $h$ obey the
same bound: $\tau_E(\eta)=\kappa j(G\eta)$, the Green operator gains
one derivative, and first-jet evaluation is continuous at these
indices. For $f_+=\mathcal Cv$, one has
\[
 d_EXf_+=\eta_0f_+,\qquad
 \mathcal RXf_+=\overline{f_+},\qquad
 \|\mathcal RXf_+\|_{H^{s+1}}
       \leq C_s\|v\|_{H^s}.
\]
The first identity uses $p_0(f_+)=0$. These bounds prove the
one-derivative gain of $K$ on both summands.

In the Hilbert coordinates, $\widehat\Pi_L$ is the ordinary spherical
spectral projection on each summand, with conjugation on the second.
It is orthogonal and converges strongly. Radial spherical projection
commutes with pointwise conjugation, so the formulas for $\mathsf N$
also give $\mathsf N\widehat\Pi_L=\widehat\Pi_L\mathsf N$.
\end{proof}

The inclusion $\mathfrak V_{s+1}\hookrightarrow\mathfrak V_s$ is
compact, as the sum of two ordinary Sobolev inclusions.
Hence $K$ and $X^2=\mathsf NK+K\mathsf N+K^2$ are compact.
The operator $X$ is not compact on this scale, because
$\mathsf N$ is an isometry on an infinite radial summand.

\begin{lemma}\label{mht:lem:staggered-quasinilpotence}
The spectrum of $X$ on $\mathfrak V_s$ is $\{0\}$.
\end{lemma}

\begin{proof}
Rotation pullback $U_\theta$ is unitary. On the positive summand,
$\mathcal RU_\theta f_+=e^{-i\theta}U_\theta\mathcal Rf_+$,
whose extra phase preserves the norm. Since $X$ raises weight by one,
\[
 U_\theta XU_\theta^{-1}=e^{i\theta}X.
\]

A nonzero spectral value would bring its entire rotation circle
into the spectrum. Polynomial spectral mapping would then put a
nonzero circle in the spectrum of the compact operator $X^2$.
The nonzero spectrum of a compact operator consists of isolated
finite-multiplicity eigenvalues, so this is impossible.
Nonemptiness of the spectrum of a bounded operator gives $\{0\}$.
\end{proof}

Fix a positive Hilbert norm on $\mathfrak g=\mathfrak{gl}_N(\mathbb C)$
and retain ordinary matrix trace in the cyclic pairing. Define
\[
 M_J=1+X\otimes\operatorname{ad}_J,\qquad
 B_J=M_J^{-1},\qquad
 B_{J,L}=(1+X_L\otimes\operatorname{ad}_J)^{-1}.
\]
The finite inverse acts on $\widehat V_L\otimes\mathfrak g$.
The tensor with $\mathfrak g$ is understood in the norms below.

\begin{theorem}\label{mht:thm:uniform-staggered-inverse}
For $s\geq3$ and every compact matrix set $\mathcal K$,
\[
 \sup_{L\geq0,\ J\in\mathcal K}
       \|B_{J,L}\|_{\mathfrak V_s\otimes\mathfrak g}<\infty.
\]
The full inverse $B_J$ is entire in the entries of $J$.
For every integer $r\geq0$ and finite parameter derivative
$\partial_J^\alpha$,
\begin{equation}\label{mht:eq:staggered-inverse-rate}
 \sup_{J\in\mathcal K}
 \|\partial_J^\alpha(B_{J,L}\widehat\Pi_L-B_J)f\|_{\mathfrak V_s}
 \leq C_{s,r,\alpha,\mathcal K}(1+L)^{-r}
                              \|f\|_{\mathfrak V_{s+r}}.
\end{equation}
Thus the inverses and all finite background responses converge
uniformly on bounded smooth sets, locally uniformly in $J$.
\end{theorem}

\begin{proof}
The spectral-radius formula gives
$\|X^n\|^{1/n}\to0$. Consequently
\[
 B_J=\sum_{n\geq0}(-X\otimes\operatorname{ad}_J)^n
\]
converges in operator norm uniformly on compact matrix sets and
defines an entire inverse.

Retain the nonsmoothing block before compressing:
\[
 \widehat X_L=\mathsf N+\widehat\Pi_LK\widehat\Pi_L.
\]
Compactness of $K$ and strong convergence of the orthogonal
projections imply $\widehat X_L\to X$ in operator norm.
Indeed, $(1-\widehat\Pi_L)K$ converges uniformly on the unit ball
because its image under $K$ is relatively compact. Apply the same
argument to the adjoint compact operator to control
$K(1-\widehat\Pi_L)$.

The operator $\widehat X_L$ commutes with $\widehat\Pi_L$.
On its range it is $X_L$, and on its complement it is $\mathsf N$.
Uniform boundedness of $B_J$ on $\mathcal K$ and the resolvent
Neumann series therefore give a uniform inverse bound for
$1+\widehat X_L\otimes\operatorname{ad}_J$ at large $L$.
Restriction gives the required bound for $B_{J,L}$.
The finitely many remaining polynomial inverses are bounded on
$\mathcal K$, so one constant works for every $L$.

For the rate set $u=B_Jf$ and $u_L=B_{J,L}\widehat\Pi_Lf$.
Projecting $M_Ju=f$ and subtracting the finite equation gives
\begin{equation}\label{mht:eq:staggered-Galerkin-error}
 u_L-u=-(1-\widehat\Pi_L)u
   +B_{J,L}\widehat\Pi_L(X\otimes\operatorname{ad}_J)
                            (1-\widehat\Pi_L)u.
\end{equation}
The adapted spectral tail obeys
\[
 \|(1-\widehat\Pi_L)v\|_{\mathfrak V_s}
     \leq C_r(1+L)^{-r}\|v\|_{\mathfrak V_{s+r}}.
\]
Apply it to $u$, using the finite inverse bound at index $s$ and
the full inverse bound at index $s+r$. This proves the rate for
$\alpha=0$.

Enlarge $\mathcal K$ by a fixed small compact complex neighborhood.
The same rate holds uniformly there. Cauchy estimates in the
finite-dimensional matrix entries give every fixed parameter
derivative in \eqref{mht:eq:staggered-inverse-rate}.
The inverses at different indices agree on smooth vectors, since
the same norm-convergent Neumann series defines them.
\end{proof}

These estimates concern the actual linear inverses on the declared
Hilbert scale. They will control smoothly smeared Gaussian
contractions and their source derivatives. A local interacting
quantum comparison still requires its cyclic counterterm domain,
normal contacts, boundary data and prepared state.
